درخت
\lr{Huffman}
را از پایین به بالا، با ادغام دو کم‌فراوانی‌ترین گره در هر گام، می‌سازیم:
\begin{flalign*}
    &A:5,\ B:9,\ C:12,\ D:13,\ E:16,\ F:45 && \\
    &(A{+}B):14,\ C:12,\ D:13,\ E:16,\ F:45 && \\
    &(C{+}D):25,\ (A{+}B):14,\ E:16,\ F:45 && \\
    &(\,(A{+}B){+}E\,):30,\ (C{+}D):25,\ F:45 && \\
    &(\,(C{+}D){+}((A{+}B){+}E)\,):55,\ F:45 && \\
    &(F + \cdots):100 &&
\end{flalign*}

با تخصیص $0$ به شاخهٔ چپ و $1$ به شاخهٔ راست، کدها چنین به دست می‌آیند:
\begin{flalign*}
    F &\to 0 \\
    C &\to 100 \\
    D &\to 101 \\
    A &\to 1100 \\
    B &\to 1101 \\
    E &\to 111 &&
\end{flalign*}

از این رو، معادل واژهٔ
$\text{\lr{AABCDFCCCF}}$:
\begin{flalign*}
    &\underbrace{1100}_{A}\,\underbrace{1100}_{A}\,\underbrace{1101}_{B}\,\underbrace{100}_{C}\,\underbrace{101}_{D}\,\underbrace{0}_{F}\,\underbrace{100}_{C}\,\underbrace{100}_{C}\,\underbrace{100}_{C}\,\underbrace{0}_{F} && \\
    &= 1100\,1100\,1101\,100\,101\,0\,100\,100\,100\,0 &&
\end{flalign*}